\documentclass[10pt,a4paper]{amsart}
\usepackage[margin=19mm]{geometry}
\usepackage{amsmath,amssymb,mathtools}
\usepackage[scr=rsfs,cal=euler]{mathalfa}
\usepackage{xfrac}
\usepackage[unicode,hidelinks,bookmarks=false]{hyperref}
\newcommand{\ie}{i.e.\ }
\newcommand{\cf}{cf.\ }
\newcommand{\resp}{respectively\ }
\newcommand{\Subsection}{\subsection}
\providecommand{\nobreakdash}{\nobreak\hbox{-}}
\setlength{\emergencystretch}{2em}
\multlinegap0pt
\usepackage{bm}
\usepackage{bbm}
\usepackage{dsfont}
\usepackage{xspace}
\usepackage{cancel}
\usepackage{mathtools}
\usepackage{amsthm}
\makeatletter
\@ifundefined{theorem}{\newtheorem{theorem}{Theorem}}{}
\@ifundefined{lemma}{\newtheorem{lemma}[theorem]{Lemma}}{}
\makeatother
\def\butterfly{\text{butterfly}}
\def\claw{\text{claw}}
\title[Cops and Robbers on Graphs with Path Constraints]{Cops and Robbers on Graphs with Path Constraints}
\author{Alexander Clow, Erin Meger}
\date{Source: arXiv:2509.10941v1 (CC BY 4.0)}
\begin{document}
\maketitle

\noindent\emph{Source and licence.} Cops and Robbers on Graphs with Path Constraints. Version arXiv:2509.10941v1 (CC BY 4.0), distributed under the Creative Commons Attribution 4.0 International licence, as recorded in the arXiv licence metadata for this exact version and shown on the abstract page https://arxiv.org/abs/2509.10941v1. Full rights notice: licenses/arxiv-2509.10941-CC-BY-4.0.txt.\par\smallskip

\noindent\textit{Editorial scope.} Counted unit: the Lemma whose source label is \texttt{Lemma: Flail claw,butterfly Lemma} in the article (source0.tex lines 339--345, source label \texttt{Lemma: Flail claw,butterfly Lemma}), delivered together with the article's own complete original proof of it (source0.tex lines 347--373, one \texttt{proof} environment of 27 lines). No mathematics is written, repaired or re-derived: statement and proof are the article's own text line for line. Required context retained uncounted: none. Every definition, display and label of those lines is retained unchanged. Omitted from this package and not redistributed: 1--338, 346--346, 374--1097. Nothing inside a retained excerpt is omitted or altered: every retained body is delivered exactly as the article's lines are written, so no omission receipt of any kind accompanies this package. Citations: the retained excerpts carry no citation command at all, so no bibliography entry of the article is transcribed and no reference is redistributed. Reference adaptation: none was needed and none was made; every reference inside the delivered unit resolves inside it, so no unresolved reference or citation is produced. Printed slips kept literal: no printed word, symbol, hypothesis or number of the counted statement or of its proof was altered. Layout and class: the article's own class is \texttt{article} with packages including authblk, amsfonts, graphicx, tabularx, array, color, comment, amsmath, amsthm, amssymb, fullpage, xcolor, tikz, listings; the wrapper is the lane's pinned \texttt{amsart} editorial wrapper at 10pt on A4, which restates the theorem-like environments the retained text needs and adds the article's own macro definitions that the retained text needs (\texttt{\textbackslash butterfly}, \texttt{\textbackslash claw}), with extra wrapper packages (bm, bbm, dsfont, xspace, cancel, mathtools). The delivered numbering is the unit's own. Assets: no retained span contains a figure environment, a graphics-inclusion command, a PDF-page inclusion, a table or a TikZ picture; no image file is copied into this package, the package declares no asset dependency and no image file is redistributed with it. Licence: version arXiv:2509.10941v1 of this article is distributed under the Creative Commons Attribution 4.0 International licence, as recorded in the arXiv licence metadata for this exact version and shown on the abstract page https://arxiv.org/abs/2509.10941v1. A full rights notice is delivered at licenses/arxiv-2509.10941-CC-BY-4.0.txt. Characters: every retained excerpt is pure ASCII, so the delivered engine, XeLaTeX with its default Latin Modern text font, reports no missing character.
\begin{lemma}\label{Lemma: Flail claw,butterfly Lemma}
    Let $G$ be a $(\claw,\butterfly)$-free graph and $H$ an induced subgraph of $G$.
    If $H$ is a $(k,t,S)$-flail
    such that $k\geq 3$, and
    \[\{ (k,j): j \in [t]\} \cap S = \emptyset,\]
    and $(i,j),(i-1,j)\in S$ for some $i$ and $j$, then for all $q \in [t]$, $(i,q)\in S$ or $(i-1,q)\in S$.
\end{lemma}

\begin{proof}
    Let $k\geq 3$, 
    let $G$ be a $(\claw,\butterfly)$-free graph, and let $H$ an induced $(k,t,S)$-flail in $G$.
    For contradiction suppose
    \[\{ (k,j): j \in [t]\} \cap S = \emptyset,\]
    and $(i,j),(i-1,j)\in S$ for some $i$ and $j$ but there exists a $q \in [t]$ such that $(i,q),(i-1,q)\notin S$.
    Then $q\neq j$ implying $t \geq 2$.

    Since $G$ is $\claw$-free, $H$ is $\claw$-free. 
    Hence, $\{ (k,j): j \in [t]\} \cap S = \emptyset$ implies that 
    vertices $\{v_1,\dots, v_t\}$ induce a clique.
    Otherwise, since $t\geq 2$, there are distinct and non-adjacent vertices $v_a,v_b$, 
    implying the vertices $\{u_k,u_{k+1},v_a,v_b\}$ induce a $\claw$, which is a contradiction.
    Suppose then that
    vertices $\{v_1,\dots, v_t\}$ induce a clique.

    Since
    vertices $\{v_1,\dots, v_t\}$ induce a clique, vertices $v_q$ and $v_j$ are adjacent.
    Trivially, $v_q$ and $v_j$ are both adjacent to $u_{k+1}$.
    By our assumption that $(i,q),(i-1,q)\notin S$, $v_q$ is not adjacent to $u_i$ or $u_{i-1}$.
    Given we assumed  $(i,j),(i-1,j)\in S$, we note that $v_j$ is adjacent to $u_i$ and $u_{i-1}$.
    Finally, we note that since $u_1\dots u_{k+1}$ is an induced path by the definition of being a $(k,t,S)$-flail,
    vertices $u_i$ and $u_{i-1}$ are adjacent, while $i < k$ since $(k,j)\notin S$, implies that both $u_i$ and $u_{i-1}$ are non-adjacent to $u_{k+1}$.
    But this implies vertices $\{u_{i-1},u_i,u_{k+1},v_j,v_q\}$ induces a $\butterfly$.
    Since $H$ is an induced subgraph of $G$, which is $\butterfly$-free, this is a contradiction.
    Thereby completing the proof.
\end{proof}

\end{document}
